\section{2022-2023学年线性代数II（H）期末考前练习答案}

\begin{enumerate}
    \item 设
    \[
        S_1=\frac{A+A^*}{2},\quad S_2=\frac{A-A^*}{2}.
    \]
    则 $S_1^*=S_1,\ S_2^*=-S_2$，且 $A=S_1+S_2$. 又
    \[
        S_1S_2-S_2S_1=\frac12(A^*A-AA^*),
    \]
    故 $A$ 正规当且仅当 $S_1S_2=S_2S_1$.

    若还有 $A=H+K$，其中 $H^*=H,\ K^*=-K$，则 $A^* = H-K$，联立可得
    \[
        H=\frac{A+A^*}{2},\quad K=\frac{A-A^*}{2},
    \]
    所以分解唯一. 实矩阵情形同理，将 $*$ 换为转置即可.

    \item 设 $A=(\alpha_1,\ldots,\alpha_n)$，$A^{\mathrm{T}}=(\beta_1,\ldots,\beta_n)$，其中 $\alpha_i,\beta_i$ 分别为 $A,A^{\mathrm{T}}$ 的列向量.

    若 $A$ 正规，则
    \[
        (\langle \alpha_i,\alpha_j\rangle)_{ij}=A^{\mathrm{T}}A=AA^{\mathrm{T}}
        =(\langle \beta_i,\beta_j\rangle)_{ij}.
    \]
    因而两组列向量的 Gram 矩阵相同，存在正交变换 $Q$ 使得 $Q\alpha_i=\beta_i$，即 $QA=A^{\mathrm{T}}$. 任意正交变换均可写成有限个镜面映像之积，故存在镜面映像 $\pi_1,\ldots,\pi_t$ 使得
    \[
        A^{\mathrm{T}}=\pi_t\cdots\pi_1A.
    \]

    反之，若 $A^{\mathrm{T}}=QA$，其中 $Q$ 是正交矩阵，则：
    \[
        AA^{\mathrm{T}}=(A^{\mathrm{T}})^{\mathrm{T}}A^{\mathrm{T}}=(QA)^{\mathrm{T}}(QA)=A^\mathrm{T}Q^\mathrm{T}QA=A^{\mathrm{T}}A,
    \]
    故 $A$ 正规.

    \item 该变换的矩阵为
    \[
        A=R_z(\theta)R_x(\varphi)=
        \begin{pmatrix}
            \cos\theta&-\sin\theta\cos\varphi&\sin\theta\sin\varphi\\
            \sin\theta&\cos\theta\cos\varphi&-\cos\theta\sin\varphi\\
            0&\sin\varphi&\cos\varphi
        \end{pmatrix}.
    \]
    它是 $\mathbf{R}^3$ 中保持定向的正交变换，故为绕某条轴旋转某角度. 任意三维旋转都可表示为两个过旋转轴的平面的镜面映像之积：若旋转角为 $\alpha$，取两个包含旋转轴且夹角为 $\alpha/2$ 的平面，其镜面映像之积即为该旋转. 因此本题所求两个镜面映像可取为上述两平面关于标准内积的反射. 若需要显式平面，先求 $A$ 关于特征值 $1$ 的单位特征向量作为旋转轴，再取任意过该轴的平面及其旋转 $\alpha/2$ 后的平面即可.

    \item
    \begin{enumerate}
        \item 特征多项式为 $\lambda^2-a$.
        若 $a>0$，则有两个不同实特征值 $\pm\sqrt{a}$，所有不变子空间为
        \[
            \{0\},\quad \spa(1,\sqrt{a})^{\mathrm{T}},\quad \spa(1,-\sqrt{a})^{\mathrm{T}},\quad \mathbf{R}^2.
        \]
        若 $a=0$，矩阵为一个二阶幂零块，所有不变子空间为
        \[
            \{0\},\quad \spa(1,0)^{\mathrm{T}},\quad \mathbf{R}^2.
        \]
        若 $a<0$，在 $\mathbf{R}$ 上无特征值，所有不变子空间为
        \[
            \{0\},\quad \mathbf{R}^2.
        \]

        \item 由 $T^n=O,\ T^{n-1}\neq O$ 知 $T$ 是一个 $n$ 阶幂零 Jordan 块. 取 Jordan 链 $v_1,\ldots,v_n$，满足
        \[
            Tv_1=0,\quad Tv_i=v_{i-1}\quad(i\geqslant 2).
        \]
        则全部不变子空间为
        \[
            \{0\},\ \spa(v_1),\ \spa(v_1,v_2),\ \ldots,\ \spa(v_1,\ldots,v_n)=V.
        \]

        \item 由计算可知，关于特征值 $0$ 的 Jordan 型为一个二阶块，关于特征值 $1$ 的 Jordan 型为一个三阶块. 可取 Jordan 链
        \[
            u_1=(0,-1,3,-3,1)^{\mathrm{T}},\quad
            u_2=(-1,3,-3,1,0)^{\mathrm{T}},
        \]
        满足 $Au_1=0,\ Au_2=u_1$；以及
        \[
            w_1=(0,0,1,-2,1)^{\mathrm{T}},\quad
            w_2=(0,0,-1,1,0)^{\mathrm{T}},\quad
            w_3=(0,0,1,0,0)^{\mathrm{T}},
        \]
        满足 $(A-I)w_1=0,\ (A-I)w_2=w_1,\ (A-I)w_3=w_2$.

        因不同特征值的广义特征子空间可直和分解，全部不变子空间为
        \[
            U_i\oplus W_j,\quad 0\leqslant i\leqslant 2,\ 0\leqslant j\leqslant 3,
        \]
        其中
        \[
            U_0=\{0\},\ U_1=\spa(u_1),\ U_2=\spa(u_1,u_2),
        \]
        \[
            W_0=\{0\},\ W_1=\spa(w_1),\ W_2=\spa(w_1,w_2),\ W_3=\spa(w_1,w_2,w_3).
        \]
    \end{enumerate}

    \item 该矩阵 $A$ 满足 $A^3=O$，且
    \[
        A^2=
        \begin{pmatrix}
            0&0&120&560\\
            0&0&0&0\\
            0&0&0&0\\
            0&0&0&0
        \end{pmatrix}\neq O.
    \]
    因此 $A$ 的 Jordan 型含一个 $J_3(0)$ 块和一个 $J_1(0)$ 块.

    若 $B^2=A$，则 $B$ 也是幂零矩阵. $4$ 阶幂零矩阵的 Jordan 块平方后，非零块大小只能由原块大小决定：$J_4(0)^2$ 的 Jordan 型为 $J_2(0)\oplus J_2(0)$，$J_3(0)^2$ 的 Jordan 型为 $J_2(0)\oplus J_1(0)$，更小块也不可能产生 $J_3(0)$. 因此 $B^2$ 不可能有 $J_3(0)$ 块，矛盾. 故不存在复矩阵 $B$ 使 $B^2=A$.

    当然也有更简单的办法：若 $B^2=A$，则 $B^6=A^3=O$，说明 $B$ 也是 $4$ 阶幂零矩阵，故 $B^4=A^2=O$，与 $A^2\neq O$ 矛盾.

    \item 设 $X=(\alpha_1,\alpha_2,\alpha_3)$，$Y=(T\alpha_1,T\alpha_2,T\alpha_3)$. 则 $T$ 在自然基下的矩阵为 $YX^{-1}$，计算得
    \[
        [T]=\frac13
        \begin{pmatrix}
            1&-2&-2\\
            -2&1&-2\\
            -2&-2&1
        \end{pmatrix}.
    \]
    该矩阵对称且平方为 $I$，特征值为 $1,1,-1$. 因此 $T$ 保持内积，是关于 $-1$ 特征空间正交补的镜面映像. 事实上 $-1$ 特征向量可取 $(1,1,1)^{\mathrm{T}}$，故 $T$ 是关于平面 $x+y+z=0$ 的镜面映像.

    \item
    \begin{enumerate}
        \item 第一个位置的加性与齐性、共轭对称性均容易验证. 且
        \[
            \langle \vec{x}, \vec{x}\rangle_V=a(x_2-x_1)^2+bx_2^2+x_3^2 \geqslant 0 \quad(a,b>0).
        \]
        若其为 $0$，故 $x_3=0,x_2=0,x_2-x_1=0$，等价于 $\vec{x}$ 为零向量，因此正定性也满足. 故 $\langle \cdot, \cdot \rangle_V$ 是内积.

        \item 对自然基作 Gram-Schmidt 正交化，计算可得
        \[
            \frac{1}{\sqrt a}(1,0,0),\quad
            \frac{1}{\sqrt b}(1,1,0),\quad
            (0,0,1)
        \]
        为一组满足要求的规范正交基.

        \item 设 $\beta=(\beta_1,\beta_2,\beta_3)$. 由
        \[
            \langle x,\beta\rangle_V
            =a(x_2-x_1)(\beta_2-\beta_1)+bx_2\beta_2+x_3\beta_3
        \]
        与 $x_1+x_2+x_3$ 比较系数，得
        \[
            \beta_2-\beta_1=-\frac1a,\quad b\beta_2=2,\quad \beta_3=1.
        \]
        故
        \[
            \beta=\left(\frac1a+\frac2b,\frac2b,1\right).
        \]
    \end{enumerate}

    \item $l_1$ 的方向向量为 $(1,-1,3)$. $l_2$ 的方向向量为
    \[
        (a,2,1)\times(1,-1,-1)=(-1,a+1,-a-2).
    \]
    它不可能与 $(1,-1,3)$ 平行，故两直线不存在平行或重合情形.

    将 $l_1$ 的参数方程代入 $l_2$ 的两个平面方程，得
    \[
        (a+1)t-2=0,\quad -t+1+d=0.
    \]
    因此两直线相交当且仅当
    \[
        (a+1)(d+1)=2.
    \]
    否则两直线异面.
\end{enumerate}

\clearpage
